\documentclass[10pt]{article}
\usepackage[a4paper,margin=0.82in]{geometry}
\usepackage[T1]{fontenc}
\usepackage{txfonts}
\usepackage[hidelinks]{hyperref}
\setlength{\parindent}{0pt}
\setlength{\parskip}{5pt}

\begin{document}

\begin{flushright}
School of Computer Science and Engineering\\
Tianjin University of Technology\\
Tianjin, China\\
May 11, 2026
\end{flushright}

Editor-in-Chief\\
\textit{Future Generation Computer Systems}

Dear Editor,

We would like to submit our manuscript entitled ``Towards Trusted Federated Learning via Adaptive Trust-Flow in Edge Computing: A Study on Key Technologies'' for consideration for publication in \textit{Future Generation Computer Systems}.

We declare that this manuscript presents original research that has not been published previously and is not currently under consideration for publication elsewhere. All authors have read and approved the submitted manuscript. We also declare that there are no known competing financial interests or personal relationships that could have appeared to influence the work reported in this paper.

In this work, we propose a trust-flow-driven trusted federated learning framework for edge computing environments. The framework integrates client admission, runtime trust monitoring, dual-stream reputation evolution, and layer-wise robust aggregation from a hardware-software co-design perspective. It aims to address the combined challenges of open edge participation, stealthy cross-round attacks, strong Non-IID data distributions, and false-positive control in trusted federated learning.

The main contributions of this manuscript are as follows. First, we design a TEE-anchored dynamic admission mechanism that combines remote attestation, runtime feature monitoring, entropy analysis, and Kalman-based trust evolution to move client admission from one-time verification toward continuous trust assessment. Second, we develop an orthogonal dual-stream reputation model that separates historical utility from instant risk, allowing the system to protect legitimate long-tail clients while rapidly responding to backdoor-like anomalies. Third, we propose a layer-wise risk-gated aggregation strategy that uses clean reference guidance, dynamic clipping, and survivor weight re-normalization to suppress malicious updates without discarding useful heterogeneous contributions. Finally, experiments on Flower with CIFAR-10, 20 clients, and six mixed attack types show that the proposed mechanism maintains an average accuracy of 92.31\% while limiting the attack success rate to 10.21\%, with the permanent false positive rate on benign clients approaching zero.

We believe that this manuscript is well aligned with the scope of \textit{Future Generation Computer Systems}, particularly its interest in edge computing, distributed intelligent systems, trustworthy computing, and secure service infrastructures. We hope that the manuscript will be suitable for consideration by the journal.

Thank you very much for your time and consideration. We look forward to receiving comments from the reviewers. If you have any questions, please do not hesitate to contact the corresponding author, Dr. Chao Bu, at \href{mailto:bc_0722@163.com}{bc\_0722@163.com}.

Yours sincerely,\\[1.2em]
Jiakai Zhou and Chao Bu

\end{document}
